\begin{proposition}[Hirzebruch germs cannot occur in a torus degeneration]\hypertarget{tag-01B5}{}\label{cdp:st:A.5.15}
\StatementRecord{cdp:st:A.5.15}{proved}{01B5}{\SeriesRef{cdp:eq:A.J.33a}, \SeriesRef{cdp:eq:A.J.33b}, \SeriesRef{cdp:eq:A.J.35}, \SeriesRef{cdp:eq:A.J.34}, \SeriesRef{cdp:eq:A.J.37}; \cite[Theorem L1 and Example 2.4(a)]{Nemethi1991Join}; \cite[Theorem 4.1 and Lemma 4.6]{MPW2005Hirzebruch}; \cite[Theorem 4.2(2)]{MPW2005Hirzebruch}}

Let $f:Y\to C_0$ be a proper holomorphic map from a
smooth complex threefold to a smooth curve, with smooth
nearby fibres complex two-tori. There is no point of
a fibre at which, in holomorphic ambient coordinates
$(x,y,z)$ and a coordinate on the base, the map has
the equation
\[
 f(x,y,z)=u(x,y,z)\bigl(x^a-y^bz^c\bigr)^m,
 \quad a\ge2,\quad b,c,m\ge1,                                    \tag{A.J.33}\hypertarget{tag-0062}{}\label{cdp:eq:A.J.33}
\]
where $u$ is a holomorphic unit. No Betti-number
assumption on $Y$ is needed. This includes the pinch
point, every germ $x^2-y^{2k}z^{2r+1}$ with
$k\ge1$, $r\ge0$, and the reducible germ
$x^2-y^2z^2$.
\end{proposition}
